\chapter*{Prefacio}
\addcontentsline{toc}{chapter}{Prefacio}

Este manual integrado parte del manual base del curso y de una idea sencilla: un estudiante no aprende diseño experimental leyendo una carpeta de documentos, sino siguiendo una secuencia de problemas que se vuelven cada vez mas claros.

El material disponible del curso contiene prácticas, presentaciones, reportes, hojas de cálculo, notas en \LaTeX{} y proyectos estudiantiles. Este documento no conserva esa estructura original como inventario. La reorganiza sobre la base del manual existente, conserva su intención académica y la expande alrededor de preguntas de aprendizaje.

\begin{idea}
Cada capítulo debe ayudar al estudiante a responder una pregunta. La práctica aparece cuando el estudiante ya tiene suficientes elementos para comprender por qué se hace, qué mide y cómo se interpreta.
\end{idea}

La versión actual integra de manera inicial los siguientes núcleos:

\begin{itemize}
  \item lentejas, germinación, miel, azúcar, café, blanco y hongos;
  \item gatos y muestreo como ejemplo de variabilidad y duplicación;
  \item ANOVA de un factor y prueba F;
  \item ANOVA de dos factores con profesor y materia;
  \item diseño factorial \(2^2\) con cúrcuma, cloro, post-it, blanco, seis repeticiones y luxómetro;
  \item tortillas y hongos como contraste entre observación interesante y diseño experimental defendible;
  \item violeta africana, café, miel, azúcar, hongos, halo café, necrosis y Taguchi \(L_4(2^3)\);
  \item procesos manuales no balanceados;
  \item ANCOVA, MANOVA y proyectos integradores.
\end{itemize}

El objetivo no es terminar el libro en una sola generación. El objetivo es dejar una base \LaTeX{} editable, reorganizable y ampliable por agentes que trabajen sobre una base de conocimiento sin perder la forma de escritura del manual original: problema, objetivo, fundamento, procedimiento, datos, interpretación y cierre prudente.

\begin{cuidado}
Este documento sustituye al borrador narrativo pequeño como base editable de integración. El manual original del curso se conserva como fuente y referencia, pero la construcción futura debe ocurrir aquí.
\end{cuidado}
